% ─────────────────────────────────────────────────────────────────────────────
%  Chương 2 — Công việc triển khai
%  Hướng dẫn của Khoa: không hạn chế số trang, trình bày theo đề mục công việc
%  (không cần theo trình tự thời gian). Đây là chương dài nhất của báo cáo.
% ─────────────────────────────────────────────────────────────────────────────
\chapter{Công việc triển khai}
\label{ch:trienkhai}

\section{Kiến trúc tổng thể}
\label{sec:kientruc}

Hệ thống tách làm hai tiến trình độc lập. Phần xử lý ảnh viết bằng Python, chạy
như một máy chủ HTTP; phần giao diện là ứng dụng React chạy trong trình duyệt và
chỉ giao tiếp với máy chủ qua REST API cùng luồng MJPEG. Cách chia này cho phép
thay giao diện mà không đụng tới phần xử lý, và ngược lại.

\begin{figure}[H]
\centering
\begin{tikzpicture}[node distance=6mm]

% ── Nguồn ──
\node[hop] (rtsp)  {Camera IP\\\scriptsize RTSP};
\node[hop, below=of rtsp] (usb) {Webcam\\\scriptsize USB};
\node[hop, below=of usb]  (tep) {Tệp video\\\scriptsize MP4};
\begin{scope}[on background layer]
  \node[khoi, fit=(rtsp)(usb)(tep), label={[font=\scriptsize\bfseries]above:Nguồn video}] (gNguon) {};
\end{scope}

% ── Luồng xử lý mỗi camera ──
\node[hopnhan, right=16mm of usb] (worker) {\textbf{CameraWorker}\\\scriptsize một luồng mỗi camera};
\node[hop, right=10mm of worker, yshift=13mm] (yolo)  {YOLOv8n\\\scriptsize phát hiện};
\node[hop, below=5mm of yolo]                 (track) {ByteTrack\\\scriptsize bám vết};
\node[hop, below=5mm of track]                (dem)   {Bộ đếm\\\scriptsize vạch / vùng};
\node[hop, below=5mm of dem]                  (la3b)  {LA-3B\\\scriptsize lọc thuộc tính};
\begin{scope}[on background layer]
  \node[khoi, fit=(yolo)(track)(dem)(la3b),
        label={[font=\scriptsize\bfseries]above:Xử lý ảnh}] (gXuLy) {};
\end{scope}

% ── Đầu ra ──
\node[hop, right=14mm of yolo]  (dem_buf) {Bộ đệm khung hình\\\scriptsize khung mới nhất};
\node[hop, below=8mm of dem_buf] (ghi)    {Ghi hình\\\scriptsize đoạn 60 giây};
\node[hop, below=8mm of ghi]     (csdl)   {SQLite (WAL)\\\scriptsize sự kiện, số đếm};

% ── Tầng phục vụ ──
\node[hopnhan, right=14mm of ghi] (api) {\textbf{FastAPI}\\\scriptsize 38 điểm cuối};
\node[hop, above=8mm of api] (mjpeg) {Luồng MJPEG\\\scriptsize \code{multipart}};
\node[hop, below=8mm of api] (react) {Giao diện React\\\scriptsize trình duyệt};

\draw[mui] (rtsp.east)  -- (worker.west);
\draw[mui] (usb.east)   -- (worker.west);
\draw[mui] (tep.east)   -- (worker.west);
\draw[mui] (worker) -- node[nhanmui, above, sloped] {khung hình} (yolo);
\draw[mui] (yolo)  -- (track);
\draw[mui] (track) -- (dem);
\draw[muimo] (track) -- node[nhanmui, right] {luồng TM} (la3b);
\draw[mui] (yolo.east)  -- node[nhanmui, above] {đã vẽ} (dem_buf.west);
\draw[mui] (dem.east)   -- (ghi.west);
\draw[mui] (dem.east)   -- (csdl.west);
\draw[mui] (dem_buf.east) -| (mjpeg.north);
\draw[mui] (csdl.east)  -- (api.west);
\draw[mui] (ghi.east)   -- (api.west);
\draw[mui] (mjpeg.south) -- (api.north);
\draw[mui] (api) -- node[nhanmui, right] {JSON} (react);
\end{tikzpicture}
\caption{Kiến trúc hệ thống. Mỗi camera chiếm một luồng xử lý riêng, kết quả đẩy
vào bộ đệm khung hình để tầng HTTP đọc mà không phải chờ.}
\label{fig:kientruc}
\end{figure}

Bên trong máy chủ, mỗi camera được cấp một luồng riêng chạy vòng lặp: đọc khung
hình, phát hiện đối tượng, bám vết, kiểm tra vượt vạch, vẽ chú thích, rồi ghi kết
quả vào bộ đệm. Tầng HTTP không bao giờ gọi trực tiếp vào camera mà chỉ đọc ảnh
JPEG mới nhất trong bộ đệm đó. Một camera bị treo hay mất mạng vì thế không làm
đứng cả API.

\section{Luồng hoạt động đầu--cuối}
\label{sec:luonghoatdong}

Các mục sau mô tả từng thành phần một cách độc lập. Mục này đi theo chiều ngược lại:
bám theo \emph{một khung hình} từ lúc rời khỏi camera đến lúc người dùng nhìn thấy kết
quả, để thấy các thành phần nối vào nhau ở đâu.

\subsection{Chín bước của một khung hình}

\begin{enumerate}[leftmargin=*, itemsep=3pt]
  \item \textbf{Đọc.} Luồng riêng của camera gọi \code{cap.read()}. Nguồn là tệp thì
        hết video sẽ phát lại từ đầu và đặt lại số đếm; nguồn là mạng thì mất tín hiệu
        sẽ chuyển sang vòng thử lại có giãn cách tăng dần từ 3 đến 60 giây.

  \item \textbf{Phát hiện.} Khung hình đưa qua YOLOv8n \emph{đúng một lần}, dù camera
        đang chạy bao nhiêu pipeline. Tham số truyền vào là hợp của mọi lớp đối tượng
        và ngưỡng tin cậy thấp nhất trong số các pipeline — chi tiết ở
        Mục~\ref{sec:chungphathien}.

  \item \textbf{Bám vết.} Mỗi pipeline lọc lại tập hộp theo lớp và ngưỡng của riêng nó,
        rồi đưa qua bộ ByteTrack \emph{riêng}. Mã định danh vết phải độc lập vì hai
        pipeline đếm hai vạch khác nhau không được chia nhau trạng thái bám.

  \item \textbf{Lọc thuộc tính} (chỉ luồng Thông minh). Các hộp chưa từng được phán
        quyết bị cắt ra, xếp hàng đưa sang LA-3B chạy trên luồng nền. Kết quả lưu vào
        bộ nhớ đệm theo mã vết nên mỗi vết chỉ tốn một lần hỏi.

  \item \textbf{Đếm.} Điểm neo bàn chân của từng vết được làm mượt qua mười khung, đối
        chiếu với vạch hoặc vùng. Vượt ra ngoài dải đệm và giữ nguyên phía mới trong ba
        khung liên tiếp thì tính một lượt.

  \item \textbf{Ghi nhận.} Lượt vượt vạch sinh một bản ghi trong bảng \code{events} kèm
        ảnh chụp thời điểm đó. Số đếm dồn được đẩy định kỳ vào bảng \code{counts} chứ
        không ghi mỗi khung, tránh làm phình cơ sở dữ liệu.

  \item \textbf{Vẽ chú thích.} Hộp giới hạn, vạch, mũi tên chiều và bảng số được vẽ lên
        một \emph{bản sao} của khung hình. Bản gốc vẫn giữ nguyên vì màn cấu hình cần
        khung hình sạch để người dùng vẽ vạch mới.

  \item \textbf{Ghi hình.} Khung đã vẽ chú thích được ghi vào tệp MP4 của đoạn 60 giây
        đang mở. Hết đoạn thì đóng tệp, mở đoạn mới và ghi bản ghi vào \code{segments}.

  \item \textbf{Phát ra ngoài.} Khung đã vẽ được đặt vào bộ đệm một ô. Tầng HTTP đọc từ
        đó và đẩy sang trình duyệt dưới dạng MJPEG.
\end{enumerate}

\subsection{Vì sao dùng bộ đệm một ô}

Bước 9 là chỗ dễ mắc lỗi nhất trong toàn bộ chuỗi. Nếu tầng HTTP đọc thẳng từ camera
thì tốc độ xử lý bị buộc vào tốc độ của trình duyệt chậm nhất đang xem — mở thêm một
tab là số đếm chậm lại theo.

Bộ đệm một ô cắt đứt ràng buộc đó: luồng camera luôn \emph{ghi đè} khung mới nhất mà
không chờ ai, còn tầng HTTP luôn đọc được khung mới nhất mà không chặn ai. Trình duyệt
chậm chỉ nhận được ít khung hơn, chứ không làm chậm việc đếm.

Đánh đổi là trình duyệt có thể bỏ qua vài khung hình. Với việc theo dõi trực tiếp thì
điều đó không quan trọng — số đếm và sự kiện vẫn tính trên toàn bộ khung hình, chỉ phần
hiển thị bị thưa đi.

\subsection{Hai đường đọc dữ liệu}

Người dùng lấy dữ liệu ra theo hai đường tách biệt, không đường nào đi qua luồng
camera:

\begin{itemize}[leftmargin=*, itemsep=2pt]
  \item \textbf{Trực tiếp} --- trình duyệt mở kết nối MJPEG tới điểm cuối
        \code{/cameras/\{id\}/stream}, giữ mở liên tục và nhận từng khung.
  \item \textbf{Xem lại} --- trình duyệt gọi các điểm cuối JSON để lấy dòng thời gian
        và danh sách sự kiện, rồi tải đoạn video MP4 qua HTTP có hỗ trợ \code{Range}.
\end{itemize}

Tách rời như vậy có nghĩa là dừng toàn bộ camera thì trang Xem lại vẫn hoạt động bình
thường trên dữ liệu đã ghi.

\section{Thu nhận và quản lý luồng video}
\label{sec:thunhan}

Lớp \code{CameraWorker} kế thừa \code{threading.Thread}. Vòng lặp ngoài cùng phụ
trách mở nguồn và kết nối lại, vòng lặp trong đọc khung hình:

\begin{lstlisting}[language=Python, caption={Vòng lặp chính của luồng camera, rút gọn từ \code{app/worker.py}}]
def run(self):
    source = _parse_source(self.camera["source"])
    while not self._stop_event.is_set():
        cap = cv2.VideoCapture(source)
        if not cap.isOpened():
            self.status = "offline"
            self._stop_event.wait(config.RECONNECT_DELAY_S)
            continue
        self._on_connected(cap)
        try:
            self._capture_loop(cap, source)
        finally:
            cap.release()
        if not self._stop_event.is_set():
            self.status = "offline"
            self._stop_event.wait(config.RECONNECT_DELAY_S)
\end{lstlisting}

Dùng \code{Event.wait()} thay cho \code{time.sleep()} là một chi tiết nhỏ nhưng
quan trọng: khi người dùng bấm tắt camera, luồng thoát ngay thay vì phải đợi hết
khoảng chờ ba giây.

Hệ thống nhận ba loại nguồn. Chuỗi toàn chữ số hiểu là chỉ số webcam, chuỗi bắt
đầu bằng \code{rtsp://} là camera IP, còn lại là đường dẫn tệp video. Tệp video
được phát lặp vô hạn để phục vụ trình diễn mà không cần camera thật.

Trước khi lưu một camera mới, giao diện gọi \code{POST /api/cameras/probe} để thử
mở nguồn và đọc một khung hình. Người dùng biết ngay đường dẫn có đúng không thay
vì tạo xong mới phát hiện camera không kết nối được.

\section{Phát hiện và bám vết đối tượng}
\label{sec:phathien}

Mô hình YOLO nạp một lần duy nhất cho toàn tiến trình. Nạp riêng cho mỗi camera
sẽ nhân bộ nhớ GPU lên nhiều lần mà tốc độ không tăng, vì các luồng vẫn phải xếp
hàng chờ cùng một thiết bị tính toán. Trạng thái bám vết thì ngược lại, phải tách
riêng theo từng camera vì mã định danh chỉ có ý nghĩa trong phạm vi một luồng
video.

Chuỗi xử lý mỗi khung hình gồm bốn bước: suy luận YOLO với ngưỡng tin cậy cấu
hình được, khử hộp trùng bằng NMS, loại hộp có diện tích dưới ngưỡng, rồi đưa qua
ByteTrack để gán mã định danh.

Ngưỡng IoU của NMS đặt ở 0{,}5 thay vì giá trị thấp hơn thường thấy. Ở cảnh đông,
hai người đứng sát nhau tạo ra hai hộp chồng lấn khá nhiều; ngưỡng quá gắt sẽ gộp
họ thành một.

\subsection{Dùng chung một lần phát hiện cho nhiều pipeline}
\label{sec:chungphathien}

Khi một camera chạy nhiều pipeline, cách làm hiển nhiên là mỗi pipeline tự gọi YOLO
trên khung hình của mình. Cách đó sai về chi phí: bốn pipeline trên một camera thành
bốn lần suy luận cho \emph{cùng một khung hình}, trong khi bản thân phép suy luận không
phụ thuộc gì vào cấu hình của pipeline.

Điều then chốt là YOLO tính toàn bộ 80 lớp COCO trong một lần lan truyền xuôi; tham số
\code{classes} chỉ lọc ở đầu ra chứ không làm mạng chạy nhẹ đi. Đo được: phát hiện sáu
lớp chỉ tốn hơn phát hiện một lớp \textbf{2{,}5\%} thời gian. Nói cách khác, thông tin
mà pipeline thứ hai cần đã nằm sẵn trong kết quả của pipeline thứ nhất.

Vì vậy phép phát hiện được tách khỏi phép bám vết:

\begin{lstlisting}[language=Python, caption={Tách phát hiện khỏi bám vết, trích từ \code{app/detector.py}}]
def detect(frame, class_ids, conf):
    """Một lần gọi YOLO cho cả camera, dùng chung cho mọi pipeline."""
    with config.inference_guard():
        result = model(frame, classes=class_ids or None, conf=conf,
                       imgsz=config.YOLO_IMGSZ, verbose=False)[0]
    return filter_min_size(apply_nms(sv.Detections.from_ultralytics(result)))

class Detector:
    def track(self, detections):
        """Mỗi pipeline lọc lại theo lớp và ngưỡng của riêng nó, rồi tự bám vết."""
\end{lstlisting}

Luồng camera gọi \code{detect()} một lần với \textbf{hợp} các lớp và ngưỡng tin cậy
\textbf{thấp nhất} của mọi pipeline đang chạy, rồi mỗi pipeline gọi \code{track()} để
lọc lại theo cấu hình riêng. Lấy hợp và ngưỡng thấp nhất là bắt buộc: thiếu một lớp
hoặc đặt ngưỡng cao hơn của pipeline nào đó thì pipeline ấy mất dữ liệu đầu vào.

Thay đổi này được kiểm chứng là \emph{không làm đổi kết quả}: chạy đối chứng 160 khung
hình với 4828 hộp giới hạn, bản cũ và bản mới cho ra kết quả giống hệt nhau ở từng hộp.
Đây là điều bắt buộc phải chứng minh khi tái cấu trúc phần tính toán --- nhanh hơn mà
đổi kết quả thì không còn là tái cấu trúc nữa.

\subsection{Nối tiếp các lần suy luận trên GPU dùng chung}
\label{sec:khoagpu}

Mỗi camera một luồng, mà tất cả dùng chung một thiết bị tính toán. Trên nền Metal của
Apple, thư viện PyTorch \emph{không an toàn khi gọi từ nhiều luồng} --- hai luồng cùng
gọi suy luận làm tiến trình chết ngay lập tức, không để lại vết trong nhật ký ứng dụng.
Sự cố này và cách truy nguyên được thuật lại ở Mục~\ref{sec:loi-sigsegv}.

Cách khắc phục là một khoá bao quanh mọi lần suy luận, chỉ bật khi thiết bị là Metal:

\begin{lstlisting}[language=Python, caption={Nối tiếp suy luận, trích từ \code{app/config.py}}]
@contextlib.contextmanager
def inference_guard():
    """Nối tiếp các lần suy luận dùng chung một GPU Metal."""
    if device() == "mps":
        with _INFERENCE_LOCK:
            yield
    else:
        yield
\end{lstlisting}

Khoá này làm các camera phải xếp hàng chờ nhau, nên tốc độ mỗi camera giảm khi số camera
tăng. Nhưng đó là giới hạn của phần cứng chứ không phải của thiết kế: một GPU chỉ tính
được một việc tại một thời điểm, khoá chỉ làm việc xếp hàng trở nên tường minh và an
toàn thay vì ngầm định và gây sập.

\subsection{Hiệu chỉnh tham số ByteTrack}
\label{sec:bytetrack}

Tham số \code{minimum\_matching\_threshold} của ByteTrack áp lên chi phí
$1 - \text{IoU}$ chứ không phải lên IoU. Đây là chỗ dễ hiểu nhầm nhất. Đặt ngưỡng
$0{,}3$ nghĩa là đòi hỏi $\text{IoU} \geq 0{,}7$ giữa hai khung hình liên tiếp
mới chấp nhận nối vết --- một yêu cầu rất khắt khe. Ở cảnh đông người, hầu như
không cặp hộp nào đạt được, kết quả là mọi vết chết ngay sau một khung hình và hệ
thống không đếm được gì.

Giá trị $0{,}9$ (tương ứng $\text{IoU} \geq 0{,}1$) khiến vết sống lâu hơn hẳn.
Đi kèm là \code{lost\_track\_buffer} đặt ở 50 khung hình để giữ vết qua các đoạn
bị che khuất ngắn.

Số liệu thực nghiệm cho năm giá trị ngưỡng nằm ở Bảng~\ref{tab:bytetrack},
Chương~\ref{ch:ketqua}.

\section{Thuật toán đếm qua vạch}
\label{sec:dem}

Đây là phần chiếm nhiều thời gian nhất của kỳ thực tập, và cũng là phần phải làm
lại nhiều lần nhất.

\subsection{Hạn chế của cách đếm sẵn có}

Thư viện Supervision cung cấp lớp \code{LineZone} đếm dựa trên tâm hộp giới hạn.
Cách này hoạt động tốt khi đối tượng đi cắt ngang tầm nhìn, nhưng hỏng trong hai
tình huống thường gặp.

Tình huống thứ nhất là người đi thẳng về phía camera. Hộp giới hạn to dần khi
người tiến lại gần, và tâm hộp dao động lên xuống theo nhịp bước chân. Mỗi lần
dao động vắt qua vạch là một lượt đếm. Hai người đi qua có thể cho kết quả bốn
lượt.

Tình huống thứ hai là ByteTrack đổi mã định danh giữa chừng do che khuất. Vết cũ
biến mất, vết mới sinh ra ở gần đó nhưng không còn lịch sử để biết nó đang ở phía
nào của vạch.

\subsection{Ba cải tiến của \code{CentroidCrossingCounter}}

\textbf{Điểm neo là bàn chân.} Thay vì tâm hộp, thuật toán lấy điểm giữa cạnh đáy
--- xấp xỉ vị trí bàn chân chạm đất. Điểm này di chuyển đơn điệu theo hướng đi và
không dao động theo nhịp bước.

\begin{lstlisting}[language=Python, caption={Chọn điểm neo, trích từ \code{app/counter.py}}]
def _anchor_pt(self, box):
    cx = (box[0] + box[2]) / 2
    cy = box[3] if self.anchor == "foot" else (box[1] + box[3]) / 2
    return np.array([cx, cy], dtype=float)
\end{lstlisting}

\textbf{Vùng đệm và độ trễ xác nhận.} Quanh vạch có một dải đệm rộng bằng 3\% chiều
cao hộp giới hạn của chính đối tượng đang xét. Đối tượng nằm trong dải này chưa
được coi là đã sang phía bên kia. Khi vượt ra ngoài dải, phía mới còn phải giữ
nguyên trong ít nhất ba khung hình liên tiếp thì mới tính là một lượt. Hai cơ chế
này cùng chặn các dao động nhỏ quanh vạch.

Vị trí dùng để xét không phải toạ độ tức thời mà là trung bình trượt của mười
khung hình gần nhất, giúp làm mượt thêm một lần nữa.

Đơn vị của dải đệm là chiều cao đối tượng chứ không phải chiều dài vạch. Cách chọn
đơn vị này là kết quả của một lỗi phải sửa về sau, trình bày ở
Mục~\ref{sec:loi-vungdem}: lấy theo chiều dài vạch thì cùng một lượt đi qua lại
được đếm hay bị bỏ tuỳ vào việc người dùng kéo vạch dài hay ngắn, điều không có
căn cứ vật lý nào. Lấy theo chiều cao đối tượng thì thang đo tự khớp với phối
cảnh: người ở xa hộp nhỏ nên dải đệm hẹp, người ở gần hộp to nên dải đệm rộng.

\textbf{Nối lại vết theo không gian.} Khi một mã định danh mới xuất hiện, thuật
toán tìm trong số các mã vừa biến mất gần đây xem có mã nào ở vị trí đủ gần
không. Nếu có, toàn bộ lịch sử của vết cũ được chuyển sang vết mới.

\begin{lstlisting}[language=Python, caption={Nối lại vết sau khi ByteTrack đổi mã định danh}]
recently_lost = {
    tid for tid, seen in self.last_seen.items()
    if tid not in current_ids
    and 0 < frame_idx - seen <= self.relink_frames
    and len(self.centroids.get(tid, ())) >= 3
}
for i, tid in enumerate(tracker_ids):
    if tid in self.centroids:
        continue
    new_pt = self._anchor_pt(boxes[i])
    best_dist, best_lost = self.relink_dist, None
    for lost in recently_lost:
        d = float(np.linalg.norm(new_pt - self.centroids[lost][-1]))
        if d < best_dist:
            best_dist, best_lost = d, lost
    if best_lost is not None:
        for store in (self.centroids, self.sides, self.cand,
                      self.cand_cnt, self.counted):
            if best_lost in store:
                store[tid] = store.pop(best_lost)
\end{lstlisting}

Thêm vào đó là khoảng nghỉ 20 khung hình sau mỗi lượt đếm cho cùng một vết, và
điều kiện hình chiếu của điểm neo phải nằm trong phạm vi đoạn vạch chứ không phải
trên phần kéo dài vô hạn của nó.

\subsection{Xác định chiều Vào và Ra}

Vạch chia mặt phẳng làm hai nửa. Dấu của khoảng cách có hướng từ điểm neo tới
đường thẳng chứa vạch cho biết đối tượng đang ở nửa nào:

\[
d = \frac{v_x \,(p_y - a_y) - v_y \,(p_x - a_x)}{\lVert v \rVert}
\]

trong đó $a$ là điểm đầu vạch, $v$ là vector chỉ phương, $p$ là điểm neo. Chuyển
từ $d < 0$ sang $d > 0$ tính là một chiều, ngược lại là chiều kia.

Quy ước nào là Vào thì phụ thuộc cách vẽ vạch, và người dùng không đoán được. Hệ
thống vẽ sẵn hai mũi tên ngay giữa vạch trên mọi khung hình, mũi tên xanh chỉ về
phía tính là Vào, mũi tên đỏ chỉ phía Ra. Nếu ngược ý, một công tắc trong giao
diện đảo lại cả số đếm lẫn hướng mũi tên.

\subsection{Đếm trong vùng đa giác}

Vạch thẳng trả lời câu hỏi "bao nhiêu lượt đi qua". Có những tình huống cần câu trả
lời khác: "hiện có bao nhiêu người trong khu vực chờ", "đã có bao nhiêu người bước vào
khu vực hạn chế". Lớp \code{PolygonZoneCounter} phụ trách phần này.

Kiểm tra một điểm có nằm trong đa giác hay không dùng \code{cv2.pointPolygonTest}, hàm
này xử lý được cả đa giác lõm. Hai nguyên tắc của bộ đếm vạch được giữ nguyên: điểm neo
vẫn là bàn chân, và trạng thái mới phải duy trì vài khung hình liên tiếp mới được xác
nhận.

Điểm neo bàn chân đặc biệt quan trọng ở đây. Người đứng ngay sát mép vùng có hộp giới
hạn phủ một phần lên vùng nhưng bàn chân vẫn ở ngoài; lấy tâm hộp sẽ tính nhầm là đã
vào, và ở biên vùng thì trạng thái sẽ nhấp nháy liên tục.

Bộ đếm giữ hai con số: số đối tượng \emph{đang} ở trong vùng, và tổng số lượt \emph{đã}
bước vào. Con số thứ hai dùng tập \code{entered} để một người ra rồi vào lại không bị
đếm thành hai lượt.

Người dùng vẽ vùng bằng cách bấm lần lượt từng đỉnh trên ảnh chụp từ camera, kéo đỉnh
để chỉnh. Không giới hạn số đỉnh, chỉ yêu cầu tối thiểu ba.

Ở giao diện, đếm qua vạch và đếm trong vùng \emph{không} là hai lựa chọn tách rời.
Người dùng chỉ chọn "Đếm đối tượng", rồi kiểu hình vẽ ở bước sau mới quyết định thuật
toán nào chạy. Với người vận hành thì đây vốn là cùng một việc; tách thành hai mục
buộc họ phải hiểu trước sự khác nhau giữa hai thuật toán mới chọn được, trong khi nhìn
khung hình rồi vẽ thì tự khắc rõ.

Khung vẽ cũng không dựng sẵn hình nào. Bản trước đặt sẵn một vạch ngang giữa khung cho
tiện, nhưng như vậy người dùng bấm "Tiếp tục" là qua luôn mà không để ý, và pipeline
chạy với vạch đặt sai chỗ. Giờ chưa vẽ thì nút đi tiếp bị khoá, khung hình phủ mờ kèm
lời nhắc.

\subsection{Đặt vạch ở đâu}

Người đi cắt ngang tầm nhìn camera cho kết quả đếm ổn định nhất. Người đi thẳng
vào hoặc ra theo trục camera là trường hợp khó, vì họ không cắt vạch một cách dứt
khoát. Với bố trí này, vạch nên đặt đúng chỗ bàn chân bước qua, thường thấp hơn
giữa khung hình.

Người đi song song với vạch không bị đếm, vì họ không đổi phía và dải đệm chặn
các trôi dạt nhỏ.

\section{Luồng Thông minh: lọc đối tượng bằng mô tả tiếng Việt}
\label{sec:thongminh}

YOLO chỉ biết 80 lớp đối tượng của bộ dữ liệu COCO. Muốn đếm riêng \emph{người mặc áo
đỏ} hay \emph{người đeo ba lô} thì bình thường phải gán nhãn và huấn luyện lại --- việc
chiếm gần hết quỹ thời gian của một kỳ thực tập.

Luồng Thông minh giải quyết bằng cách ghép hai mô hình. YOLO vẫn phụ trách tìm người,
còn LocateAnything-3B (LA-3B) --- một mô hình thị giác--ngôn ngữ --- nhận từng khung cắt
kèm mô tả và trả lời khung đó có khớp mô tả hay không. Không cần huấn luyện gì thêm.

\subsection{Xử lý câu lệnh tiếng Việt}

Người dùng gõ cả câu, không chỉ thuộc tính. Hàm \code{parse\_command} tách câu thành hai
phần: thuộc tính cần lọc và chiều cần đếm.

\begin{table}[H]
\centering
\caption{Kết quả phân tích câu lệnh và dịch sang truy vấn cho LA-3B.}
\label{tab:parsecommand}
\small
\begin{tabularx}{\textwidth}{@{}XXll@{}}
\toprule
\textbf{Người dùng gõ} & \textbf{Thuộc tính tách được} &
\textbf{Truy vấn LA-3B} & \textbf{Chiều đếm} \\
\midrule
người đeo ba lô & người đeo ba lô & \code{person with backpack} & cả hai \\
đếm người mặc áo đỏ đi vào & người mặc áo đỏ & \code{person in red shirt} & chỉ Vào \\
người cầm ô đi ra & người cầm ô & \code{person holding umbrella} & chỉ Ra \\
người & --- & --- & bỏ qua LA-3B \\
\bottomrule
\end{tabularx}
\end{table}

Bước dịch sang tiếng Anh là bắt buộc. LA-3B huấn luyện chủ yếu trên tiếng Anh và tiếng
Trung nên không hiểu tiếng Việt: đưa thẳng câu tiếng Việt vào thì mô hình giữ lại gần
như mọi khung cắt, tức là không lọc gì cả.

Hệ thống dịch bằng từ điển chốt cứng cho các vật thể hay gặp, phần còn lại giao cho mô
hình Opus-MT. Từ điển cần thiết vì mô hình dịch sai một số từ ngắn --- chẳng hạn "ô"
thành \en{pane} thay vì \en{umbrella}. Bản dịch được nhớ lại nên mỗi câu lệnh chỉ dịch
đúng một lần.

Câu lệnh chỉ là "người" sẽ được nhận ra và bỏ qua LA-3B hoàn toàn, vì YOLO đã làm được
việc đó nhanh hơn nhiều.

\subsection{Bài toán tốc độ}

LA-3B mất khoảng một giây cho mỗi khung cắt, trong khi vòng lặp camera chạy 15 khung
hình mỗi giây. Gọi trực tiếp trong vòng lặp sẽ làm luồng video đứng hình hoàn toàn.

Ba cơ chế được dùng cùng lúc:

\textbf{Hỏi một lần cho mỗi vết.} Phán quyết lưu theo mã định danh. Một người xuất hiện
trong 300 khung hình chỉ tốn đúng một lần gọi LA-3B. Đây là lý do luồng Thông minh dùng
bộ tham số ByteTrack riêng, giữ vết lâu hơn: mỗi lần đổi mã định danh là một lần phải
hỏi lại.

\textbf{Gọi bất đồng bộ.} Một luồng phụ xử lý hàng đợi khung cắt; vòng lặp chính đẩy
khung cắt vào hàng đợi rồi đi tiếp ngay.

\textbf{Đệm lượt vượt vạch.} Đây là chi tiết dễ bỏ sót nhất. Trong khoảng một giây chờ
phán quyết, vết vẫn có thể đi qua vạch. Nếu chỉ đơn giản ẩn vết chưa có phán quyết khỏi
bộ đếm thì mọi lượt vượt vạch xảy ra trong giây đầu tiên của mỗi vết đều bị mất --- mà
ở cảnh người đi nhanh, đó là phần lớn các lượt.

\begin{lstlisting}[language=Python, caption={Ghi tạm lượt vượt vạch của vết đang chờ phán quyết, trích từ \code{app/smart\_filter.py}}]
if tid in self.pending:
    side = self._side_of(boxes[index])
    previous = self.prev_side.get(tid)
    if previous is not None and previous != side:
        delta = 1 if side > 0 else -1
        if self.flip:
            delta = -delta
        self.pending_crossings[tid] += delta
    self.prev_side[tid] = side
\end{lstlisting}

Khi luồng phụ trả về kết quả khớp, các lượt đã đệm được cộng bù vào bộ đếm; nếu không
khớp thì bỏ đi.

\subsection{Nạp mô hình}

Repo gốc của LA-3B viết cho một phiên bản thư viện cũ hơn nên phải vá ba chỗ trước khi
nạp được: bổ sung vài khoá thiếu trong \code{config.json}, trỏ tạm kiến trúc
\code{qwen3} sang \code{qwen2} vì checkpoint thực chất là Qwen2, và tạo module rỗng cho
\code{decord} với \code{lmdb} --- hai thư viện chỉ dùng cho nhánh xử lý video mà hệ
thống này không đụng tới.

Việc nạp là lazy, chỉ chạy khi có pipeline Thông minh được kích hoạt. Nếu không tìm thấy
mô hình, giao diện khoá chế độ Thông minh lại kèm hướng dẫn, phần còn lại vẫn hoạt động
bình thường.

\section{Tham số vận hành của pipeline}
\label{sec:thamso}

Ngoài hình vẽ và lớp đối tượng, mỗi pipeline còn có một nhóm tham số cho người vận
hành tự chỉnh mà không cần đụng tới mã nguồn.

\textbf{Hướng đếm} giới hạn chiều được ghi nhận. Lọc ngay trong bộ đếm chứ không lọc
danh sách sự kiện trả về; nếu lọc ở khâu sau thì con số hiển thị trên màn hình vẫn
cộng cả chiều người dùng không yêu cầu.

\textbf{Ngưỡng cảnh báo} sinh một sự kiện khi số đối tượng trong khung, hoặc trong
vùng, vượt mức đặt. Có khoảng nghỉ 60 giây giữa hai lần báo --- thiếu nó thì một đám
đông đứng yên sẽ sinh ra một cảnh báo mỗi khung hình.

\textbf{Tự động đặt lại} đưa số đếm về 0 theo chu kỳ giờ hoặc ngày. Hệ thống so sánh
mốc thời gian hiện tại với mốc lần trước; lần đầu chỉ ghi nhận mốc chứ không đặt lại,
tránh xoá sạch số đếm ngay khi vừa bật pipeline.

\textbf{Tốc độ xử lý} gói ba mức đánh đổi giữa tài nguyên máy và độ nhạy, ứng với các
cặp nhịp xử lý và ngưỡng tin cậy khác nhau. Người dùng chọn theo mô tả bằng lời;
những ai cần kiểm soát chi tiết thì mở mục tuỳ chỉnh nâng cao để kéo từng thanh
trượt.

\textbf{Lịch chạy} giới hạn AI chỉ hoạt động trong những khung giờ khai báo trước, ví
dụ chỉ đếm khách trong giờ mở cửa. Ngoài khung giờ, camera vẫn phát hình và vẫn ghi
hình --- chỉ phần nhận diện tạm nghỉ. Cách này vừa tiết kiệm tài nguyên vừa không tạo
ra số liệu rác lúc cửa hàng đóng, mà vẫn giữ được dữ liệu để xem lại nếu có việc.

Ngày trong tuần đánh số theo cách gọi của người Việt, thứ Hai là 2 cho tới Chủ nhật là
8, thay vì 0--6 như Python. Dữ liệu lưu xuống nhờ vậy đọc lên khớp luôn với nhãn hiển
thị, khỏi phải quy đổi ở hai đầu.

Trường hợp cần chú ý là khung giờ vắt qua nửa đêm, chẳng hạn 22:00 đến 06:00. Phần sau
nửa đêm thuộc về ngày hôm sau nên khi kiểm tra phải đối chiếu với danh sách ngày của
\emph{hôm trước}; nếu chỉ so với ngày hiện tại thì ca đêm sẽ bị cắt đúng lúc 0 giờ.

\subsection{Một lỗi đo tốc độ}

Khi kiểm thử phần tham số này, đặt nhịp 10 khung hình mỗi giây nhưng màn hình vẫn báo
26. Nguyên nhân nằm ở cách đo: đoạn mã cũ tính tốc độ bằng nghịch đảo \emph{thời gian
xử lý một khung hình}, tức là đo năng lực tối đa của máy chứ không phải nhịp chạy
thật, vì phần thời gian vòng lặp chủ động nghỉ không được tính vào.

Chừng nào nhịp mục tiêu còn là một hằng số chôn trong mã thì sai lệch này không ai
thấy. Nhưng khi người dùng tự đặt nhịp, con số hiển thị mâu thuẫn trực tiếp với con số
họ vừa nhập.

Cách sửa là đo khoảng cách giữa hai lần bắt đầu vòng lặp liên tiếp thay vì đo thời
gian xử lý. Thời gian xử lý vẫn được giữ lại nhưng báo cáo riêng dưới tên độ trễ ---
nó trả lời câu hỏi khác: máy có kịp xử lý không. Sau khi sửa, đặt 10 thì đo được 9{,}6.

\section{Ghi hình phân đoạn và xem lại}
\label{sec:ghihinh}

Video ghi thành nhiều tệp ngắn, mặc định mỗi tệp 60 giây, đặt trong thư mục theo
camera và theo ngày. Chia nhỏ có ba cái lợi: xoá dữ liệu quá hạn chỉ là xoá tệp,
trình duyệt tải nhanh vì tệp nhẹ, và mất điện giữa chừng chỉ hỏng đoạn đang ghi.

Nhịp ghi giảm xuống 10 khung hình mỗi giây và chiều rộng tối đa 854 điểm ảnh.
Ghi nguyên nhịp và nguyên độ phân giải sẽ ngốn ổ đĩa rất nhanh mà không giúp gì
cho việc xem lại.

Về bộ mã hoá, \code{avc1} (H.264) phát được trực tiếp trong thẻ \code{<video>},
nhưng bản OpenCV cài qua \code{pip} không phải lúc nào cũng kèm bộ mã hoá này do
vướng bản quyền. Chương trình thử \code{avc1} trước, không được thì lùi về
\code{mp4v} và ghi một dòng cảnh báo vào nhật ký.

\subsection{Hỗ trợ tua video}

Muốn thanh tua của trình duyệt hoạt động, máy chủ phải hiểu tiêu đề HTTP
\code{Range} và trả về mã 206 kèm đúng đoạn byte được yêu cầu. Thiếu phần này,
trình duyệt buộc phải tải hết tệp mới phát được và thanh tua bị vô hiệu hoá:

\begin{lstlisting}[language=Python, caption={Trả về một phần tệp theo yêu cầu Range, trích từ \code{app/routers/playback.py}}]
return StreamingResponse(
    iter_chunk(),
    status_code=206,
    media_type=media_type,
    headers={
        "Content-Range": f"bytes {start}-{end}/{file_size}",
        "Accept-Ranges": "bytes",
        "Content-Length": str(end - start + 1),
    },
)
\end{lstlisting}

Bấm vào một sự kiện trong danh sách sẽ gọi \code{GET /api/playback/at}, trả về
đoạn video chứa mốc thời gian đó cùng số giây cần tua bên trong đoạn.

\section{Tìm kiếm bằng câu tiếng Việt}
\label{sec:timkiem}

Người vận hành cần tra cứu kiểu \emph{"người đi vào hôm qua buổi sáng"} mà không
phải điền vào bốn ô lọc riêng biệt.

Hệ thống phân tích câu bằng luật thay vì gọi mô hình ngôn ngữ. Phạm vi truy vấn ở
đây hẹp và đóng: chỉ có bốn tiêu chí là đối tượng, chiều đi, camera và khoảng
thời gian. Một bộ luật vài chục dòng xử lý được toàn bộ, cho kết quả xác định,
chạy hoàn toàn ngoại tuyến, độ trễ dưới một mili giây, và không phát sinh chi phí
gọi dịch vụ bên ngoài.

Bước đầu tiên là bỏ dấu tiếng Việt để việc so khớp không phụ thuộc cách gõ:

\begin{lstlisting}[language=Python, caption={Chuẩn hoá chuỗi trước khi so khớp từ khoá}]
def strip_accents(text: str) -> str:
    text = text.replace("d", "d").replace("D", "D")
    nfkd = unicodedata.normalize("NFD", text)
    return "".join(c for c in nfkd if unicodedata.category(c) != "Mn").lower()
\end{lstlisting}

Sau đó lần lượt dò từ khoá đối tượng, từ khoá chiều đi, và cụm chỉ thời gian.
Phần chữ còn lại được đem so với tên và vị trí các camera đã khai báo. Kết quả
phân tích hiển thị lại cho người dùng dưới dạng danh sách tiêu chí, để họ biết hệ
thống đã hiểu câu hỏi ra sao và sửa lại nếu sai.

Cụm chỉ thời gian nhận được nhiều dạng: hôm nay, hôm qua, tuần này, tuần trước,
\emph{N} ngày trước, các buổi trong ngày, và khoảng giờ tường minh kiểu
\emph{"từ 8h đến 11h"}.

\section{Giao diện người dùng}
\label{sec:giaodien}

Giao diện chia thành ba trang tương ứng ba nhóm công việc.

\textbf{Giám sát} hiển thị lưới camera với ba mức 1$\times$1, 2$\times$2 và
3$\times$3, chiếm toàn bộ chiều ngang màn hình. Việc khai báo camera mới cũng đặt
ngay tại trang này thay vì tách ra một trang quản trị riêng, vì đây là thao tác
người vận hành làm thường xuyên nhất.

Cảnh báo không hiển thị thành một thanh cố định. Bản đầu tiên làm như vậy, nhưng
danh sách gộp cảnh báo của mọi camera trôi quá nhanh --- ở cảnh đông, mỗi phút vài
chục sự kiện --- và nhìn vào không biết sự kiện thuộc camera nào. Cách làm hiện tại:
bấm vào một camera đang chạy pipeline thì mở màn hình chi tiết kèm bảng cảnh báo của
đúng camera đó. Mỗi cảnh báo đi cùng ảnh chụp tại thời điểm xảy ra và hai nút xác
nhận xử lý, đóng sự kiện.

Vùng bấm để mở là toàn bộ khung hình chứ không phải một nút nhỏ ở góc. Bản trước
đặt nút phóng to ở góc trên và chỉ hiện khi rê chuột vào, người dùng thử không tìm
ra. Màn hình chi tiết có thêm chế độ toàn màn hình: hệ thống vừa gọi Fullscreen API
của trình duyệt, vừa tự phủ kín cửa sổ bằng CSS. Cần cả hai vì Fullscreen API chỉ
chạy khi có thao tác thật của người dùng và bị chặn trong một số môi trường nhúng ---
lớp CSS đảm bảo nút luôn có tác dụng, lớp API chỉ là phần thêm để ẩn nốt thanh địa
chỉ.

Lưới cũng tự co số cột theo số camera đang hiển thị. Hai camera trong lưới
2$\times$2 sẽ xếp thành hai cột chứ không để hai ô rỗng, và các hàng chia đều chiều
cao còn lại nên khung hình giãn ra lấp hết chỗ trống thay vì co về tỉ lệ 16:9.

\textbf{Cấu hình} liệt kê các pipeline đã tạo dưới dạng bảng, và chứa trình hướng
dẫn bốn bước: chọn camera, chọn bài toán, chọn luồng xử lý cùng đối tượng và vẽ
hình, cuối cùng là kiểm tra rồi kích hoạt. Trong lúc đang tạo, bảng pipeline được
ẩn đi để người dùng tập trung vào từng bước.

Bước thứ ba đặt khung vẽ bên trái và các lựa chọn bên phải trên cùng một màn hình.
Ban đầu hai phần này tách thành hai bước riêng, nhưng khi dùng thử thì thấy người
ta hay nhìn khung hình rồi mới quyết định nên đếm theo vạch hay theo vùng --- tách
ra buộc phải đi tới đi lui giữa hai bước. Ngay trên khung vẽ có công tắc chuyển
giữa \emph{vẽ vạch} và \emph{vẽ vùng}; đổi công tắc cũng là đổi luôn bài toán, vì
hai thứ này gắn liền với nhau.

Một quyết định về ngôn ngữ hiển thị: giao diện \emph{không} nhắc tên mô hình. Người
vận hành chỉ thấy hai lựa chọn "Tiêu chuẩn" và "Thông minh" cùng mô tả việc mỗi chế
độ làm được gì. Tên mô hình chỉ xuất hiện trong nhật ký hệ thống và tài liệu kỹ
thuật. Lý do là người dùng cuối không có cơ sở để chọn giữa hai cái tên mô hình,
nhưng chọn được giữa "đếm tất cả mọi người" và "chỉ đếm người mặc áo đỏ".

Toạ độ vạch lưu ở dạng chuẩn hoá trong đoạn $[0,1]$ chứ không phải điểm ảnh. Cùng
một cấu hình vì thế dùng được cho cả ảnh hiển thị đã thu nhỏ lẫn khung hình gốc
mà bộ đếm làm việc, và vẫn đúng khi camera đổi độ phân giải.

\textbf{Xem lại và Tìm kiếm} gồm trình phát video, thanh thời gian 24 giờ đánh
dấu các vùng đã ghi hình cùng vị trí sự kiện, và ô tìm kiếm tiếng Việt. Kết quả tìm
kiếm trả về dưới dạng lưới ảnh chứ không phải danh sách chữ: mỗi ô là ảnh chụp tại
đúng thời điểm sự kiện, nhìn là biết ngay có đúng thứ mình tìm hay không.

Ảnh trong lưới dùng bản thu nhỏ sinh phía máy chủ. Ảnh gốc từ camera 4K nặng khoảng
520~KB, một lưới sáu mươi kết quả sẽ phải tải hơn ba mươi megabyte; bản thu nhỏ chỉ
khoảng 10~KB, giảm gần năm mươi lần.

Các chức năng quản trị (thời hạn lưu trữ, dung lượng đĩa, nhật ký hoạt động) vẫn có
đầy đủ ở tầng API và dùng được qua tài liệu OpenAPI, nhưng không có trang giao diện
riêng. Việc dọn dữ liệu quá hạn chạy tự động theo lịch, không cần thao tác thủ công.

\section{Cơ sở dữ liệu và chính sách lưu trữ}
\label{sec:csdl}

Lược đồ gồm bảy bảng: \code{cameras}, \code{pipelines}, \code{events},
\code{counts}, \code{segments}, \code{logs} và \code{settings}. Chi tiết từng cột
xem Phụ lục~\ref{app:schema}.

\begin{figure}[H]
\centering
\begin{tikzpicture}[node distance=9mm]

\node[bang] (cameras) {\textbf{cameras}
  \nodepart{two}\ \code{id} PK\\\ name, location\\\ source, kind\\\ enabled\ };

\node[bang, right=22mm of cameras] (pipelines) {\textbf{pipelines}
  \nodepart{two}\ \code{id} PK\\\ \code{camera\_id} FK\\\ task, mode, classes\\\ line, zone, conf\\\ schedule\ };

\node[bang, below left=12mm and -6mm of pipelines] (events) {\textbf{events}
  \nodepart{two}\ \code{id} PK\\\ \code{camera\_id} FK\\\ \code{pipeline\_id} FK\\\ ts, kind, label\\\ snapshot\ };

\node[bang, below right=12mm and -6mm of pipelines] (counts) {\textbf{counts}
  \nodepart{two}\ \code{id} PK\\\ \code{pipeline\_id} FK\\\ ts, in, out\\\ occupancy\ };

\node[bang, below=26mm of cameras] (segments) {\textbf{segments}
  \nodepart{two}\ \code{id} PK\\\ \code{camera\_id} FK\\\ start\_ts, end\_ts\\\ path, duration\ };

\node[bang, right=16mm of counts] (settings) {\textbf{settings}
  \nodepart{two}\ \code{key} PK\\\ value\ };

\draw[mui] (cameras.east)  -- node[nhanmui, above] {1..n} (pipelines.west);
\draw[mui] (cameras.south) -- node[nhanmui, left]  {1..n} (segments.north);
\draw[mui] (pipelines.south) -- (events.north);
\draw[mui] (pipelines.south) -- (counts.north);
\draw[mui] (cameras.south east) -- (events.north west);
\end{tikzpicture}
\caption{Lược đồ cơ sở dữ liệu. Xoá một camera sẽ kéo theo các pipeline gắn với
nó nhờ ràng buộc khoá ngoại.}
\label{fig:csdl}
\end{figure}

SQLite bật chế độ WAL để nhiều luồng camera ghi đồng thời mà không khoá lẫn nhau.
Một khoá ở tầng Python bao quanh mọi thao tác ghi, tránh tình trạng
\emph{database is locked} khi số camera tăng.

Dữ liệu quá hạn được dọn tự động sáu giờ một lần. Bộ dọn xoá cả bản ghi trong cơ
sở dữ liệu lẫn tệp ảnh và tệp video tương ứng; xoá thiếu một trong hai sẽ để lại
tệp mồ côi chiếm ổ đĩa hoặc bản ghi trỏ tới tệp không còn tồn tại.
